\documentclass[11pt,oneside]{book}
\usepackage{springer_monograph}

\hypersetup{
    pdftitle={Chapter 22: Covariates, Longitudinal Profiling, and Repeated Measures in Mass Spectrometry},
    pdfauthor={Lucas VHH TRAN}
}

\begin{document}

\setcounter{chapter}{21} % Sets chapter counter so next is Chapter 22
\chapter{Covariates, Longitudinal Profiling, and Repeated Measures in Mass Spectrometry: Linear Mixed Models, Variance Partitioning, and Longitudinal Splines}
\label{ch:covariates_repeated}

\noindent
\textbf{Lucas VHH TRAN}\footnote{CK-CARE Davos, Christine K\"uhne -- Center for Allergy Research and Education, Davos, Switzerland. Email: \texttt{lucas.tran@ck-care.ch}}\\
\rule{\textwidth}{0.8pt}

\begin{abstract}
\noindent
Quantitative mass spectrometry (MS) experiments in clinical and translational cohorts routinely depart from simple, independent two-group cross-sectional designs. Modern biomolecular investigations frequently feature paired pre- and post-treatment biopsies, multi-day technical batch runs, complex factorial interventions, and multi-timepoint longitudinal monitoring of the same subjects. Applying standard Ordinary Least Squares (OLS) models, two-sample Student's $t$-tests, or standard unblocked ANOVA to such structured datasets violates the core assumption of independent and identically distributed (i.i.d.) observations. This analytical error---known as \textit{pseudoreplication}---underestimates parameter standard errors, produces artificially deflated $p$-values, and inundates the literature with irreproducible biomarker candidates. Conversely, failing to adjust for technical confounders (such as LC-MS run day, operator, or autosampler tray) inflates residual variance and obscures genuine biological effects. 

This chapter establishes a rigorous mathematical, methodological, and computational framework for modeling complex experimental structure in high-dimensional proteomics and metabolomics. We develop the matrix formulation of Linear Mixed-Effects Models (LMMs), derive Henderson's mixed-model equations, and contrast Restricted Maximum Likelihood (REML) against Maximum Likelihood (ML) estimation. We formalize the Intraclass Correlation Coefficient (ICC) and provide an analytical proof of the variance inflation factor induced by ignoring repeated measures. We evaluate alternative covariance structures---including Compound Symmetry, First-Order Autoregressive ($\operatorname{AR}(1)$), and Random Coefficient models---and establish natural cubic splines for capturing non-linear pharmacodynamic time courses. Furthermore, we bridge linear mixed models with empirical Bayes moderation through the \texttt{dream} (\textit{Differential expression for repeated measures}) framework. The chapter features a tutor diagnostic heuristics checklist, an end-to-end reproducible R laboratory, a critical review of the ``Five Deadly Sins'' of longitudinal MS analysis, and multi-tier exercises with exhaustive, step-by-step worked solutions.
\end{abstract}

\noindent\textbf{Keywords:} Linear Mixed-Effects Models, Repeated Measures, Longitudinal Profiling, Intraclass Correlation Coefficient (ICC), REML, Pseudoreplication, Splines, Empirical Bayes Moderation, \texttt{dream}, \texttt{variancePartition}, Kenward-Roger.

\vspace{1.5em}

% =============================================================================
\section{Tutor's Methodological Roadmap for Longitudinal and Structured Omics}
% =============================================================================

Before writing computational scripts or specifying design formulas, the quantitative scientist must map the dependencies inherent to the experimental design. Table~\ref{tab:repeated_roadmap} outlines the four sequential checkpoints required to analyze complex mass spectrometry cohorts.

\begin{table}[ht]
\centering
\caption{Tutor's Methodological Roadmap: Four Checkpoints for Complex and Longitudinal Omics}
\label{tab:repeated_roadmap}
\small
\begin{tabularx}{\textwidth}{l p{3.8cm} p{4.2cm} X}
\toprule
\textbf{Stage} & \textbf{Physical Mechanism} & \textbf{Key Analytical Objective} & \textbf{Primary Failure Mode \& Remedy} \\
\midrule
\textbf{1. Dependency Mapping} & Repeated sampling of identical patients across timepoints or tissue biopsies \cite{hurlbert1984pseudoreplication, diggle2002analysis}. & Distinguish between-subject biological factors from within-subject longitudinal kinetics. & Pseudoreplication: treating repeated draws as independent. Remedy: Linear mixed model with random intercepts \cite{laird1982random}. \\
\addlinespace
\textbf{2. Variance Partitioning} & Total variation split into patient baseline, treatment response, and technical noise \cite{hoffman2016variancepartition}. & Quantify Intraclass Correlation Coefficient (ICC); audit technical vs. biological variance. & Technical batch swamping biological signal. Remedy: Model batch as fixed or random blocking factor \cite{choi2014msstats}. \\
\addlinespace
\textbf{3. Trajectory Selection} & Non-linear pharmacodynamic time curves or decaying biomarker kinetics \cite{hastie1990generalized, wood2017generalized}. & Parameterize longitudinal dynamics via natural cubic splines or polynomial terms. & Forcing linear fit onto peaked/transient responses. Remedy: Natural spline basis matrix \texttt{ns(time, df=3)}. \\
\addlinespace
\textbf{4. Moderated Inference} & High-dimensional small-sample testing across thousands of proteins \cite{hoffman2021dream, smyth2004linear}. & Shrink feature residual variances; approximate degrees of freedom via Kenward-Roger \cite{kenward1997small}. & Deflated standard errors in unbalanced cohorts. Remedy: Deploy \texttt{dream} with REML moderation \cite{kuznetsova2017lmertest}. \\
\bottomrule
\end{tabularx}
\end{table}

\begin{notebox}
\textbf{Tutor's Methodological Advice: The True Cost of Pseudoreplication}\\
Treating repeated measures from the same patient as independent samples violates the Gauss-Markov theorem \cite{hurlbert1984pseudoreplication}. If $m$ samples per patient exhibit an intraclass correlation $\rho = 0.60$, the effective sample size is not $N \times m$, but rather:
$$N_{\text{eff}} = \frac{N \times m}{1 + (m - 1)\rho}$$
For 10 patients sampled 4 times ($n_{\text{total}} = 40$), the effective sample size is only $40 / (1 + 3 \times 0.60) = 40 / 2.8 \approx 14.3$! Running an unblocked $t$-test invents 25 degrees of freedom out of thin air, generating massive false-positive discovery rates.
\end{notebox}

% =============================================================================
\section{Introduction: The Structural Complexity of Omics Experiments}
% =============================================================================

A pervasive misconception in high-throughput biomolecular analysis is that differential abundance analysis is fundamentally an exercise in two-group hypothesis testing. In clinical proteomics, drug discovery, and biomarker validation, samples are almost never purely independent draws from an exchangeable population. Real-world mass spectrometry studies are characterized by multiple layers of hierarchy, grouping, and repeated sampling:
\begin{itemize}
    \item \textbf{Longitudinal Time Courses:} In pharmacodynamic and vaccination studies, biospecimens (e.g., plasma, serum, tissue biopsies) are collected from the same individual patient across multiple consecutive timepoints (e.g., baseline, 2 hours, 24 hours, 7 days post-treatment) \cite{fitzmaurice2012applied}.
    \item \textbf{Matched and Paired Cohorts:} Studies often utilize paired designs where diseased and adjacent non-diseased tissue are resected from the same cancer patient, or where subjects serve as their own internal controls in cross-over trials.
    \item \textbf{Hierarchical / Nested Sampling:} High-content single-cell or spatial MS workflows measure multiple single cells, laser-capture microdissection (LCM) spots, or technical instrument replicates nested within specific tissue blocks, patients, or microtiter plates \cite{gatto2023qfeatures}.
    \item \textbf{Extraneous Technical Confounders:} Samples analyzed over weeks or months inevitably span multiple LC columns, chromatography gradients, instrument calibrations, digestion batches, and operator shifts \cite{leek2010tackling, johnson2007adjusting}.
\end{itemize}

When structured or longitudinal observations are analyzed with naive methods that treat each mass spectrum as an independent realization, standard errors are artificially compressed by a factor that depends on the magnitude of the within-subject correlation. This structural mismatch produces severe false-positive inflation. Conversely, when large patient-to-patient baseline variability is left unmodeled in paired designs, genuine drug responses are swallowed by inter-individual heterogeneity, producing false negatives. The statistical model must faithfully reflect the physical architecture of the experiment.

% =============================================================================
\section{Mathematical Foundations of Linear Mixed-Effects Models}
% =============================================================================

\subsection{Matrix Formulation and Hierarchical Structure}

To accommodate both population-level trends (fixed effects) and subject-specific deviations (random effects), we formalize the Linear Mixed-Effects Model (LMM). Let $N$ denote the total number of subjects (clusters), with subject $i \in \{1, \dots, N\}$ measured at $n_i$ occasions, yielding a total sample size of $n_{\text{total}} = \sum_{i=1}^N n_i$. For a given mass spectrometry feature $g$ (peptide, protein, or metabolite), let $\mathbf{y}_i \in \mathbb{R}^{n_i}$ denote the vector of observed $\log_2$-transformed abundance values for subject $i$. The general mixed-effects model is formulated as \cite{pinheiro2000mixed, laird1982random}:
\begin{equation}
\label{eq:lmm_subject}
\mathbf{y}_i = \mathbf{X}_i \boldsymbol{\beta} + \mathbf{Z}_i \mathbf{u}_i + \boldsymbol{\varepsilon}_i, \quad i = 1, \dots, N
\end{equation}
where:
\begin{itemize}
    \item $\mathbf{X}_i \in \mathbb{R}^{n_i \times p}$ is the design matrix for the $p$-dimensional vector of fixed-effect regression coefficients $\boldsymbol{\beta} \in \mathbb{R}^p$ (e.g., treatment group, sex, age, time).
    \item $\mathbf{Z}_i \in \mathbb{R}^{n_i \times q}$ is the design matrix for the $q$-dimensional vector of subject-specific random effects $\mathbf{u}_i \in \mathbb{R}^q$ (e.g., random intercepts, random time slopes).
    \item $\mathbf{u}_i \sim \mathcal{N}_q(\mathbf{0}, \mathbf{G})$ represents the random effects vector, assumed independent across subjects, where $\mathbf{G} \in \mathbb{R}^{q \times q}$ is a symmetric positive semi-definite covariance matrix.
    \item $\boldsymbol{\varepsilon}_i \sim \mathcal{N}_{n_i}(\mathbf{0}, \mathbf{R}_i)$ is the unobserved residual error vector, with residual covariance matrix $\mathbf{R}_i \in \mathbb{R}^{n_i \times n_i}$. In standard models, $\mathbf{R}_i = \sigma_e^2 \mathbf{I}_{n_i}$.
    \item The random effects $\mathbf{u}_i$ and residual errors $\boldsymbol{\varepsilon}_i$ are mutually independent: $\operatorname{Cov}(\mathbf{u}_i, \boldsymbol{\varepsilon}_i) = \mathbf{0}$.
\end{itemize}

Stacking across all $N$ subjects, the complete system is expressed compactly as:
\begin{equation}
\label{eq:lmm_full}
\mathbf{y} = \mathbf{X} \boldsymbol{\beta} + \mathbf{Z} \mathbf{u} + \boldsymbol{\varepsilon}
\end{equation}
where $\mathbf{y} \in \mathbb{R}^{n_{\text{total}}}$, $\mathbf{X} \in \mathbb{R}^{n_{\text{total}} \times p}$, $\mathbf{Z} = \operatorname{diag}(\mathbf{Z}_1, \dots, \mathbf{Z}_N) \in \mathbb{R}^{n_{\text{total}} \times Nq}$, $\mathbf{u} \sim \mathcal{N}(\mathbf{0}, \widetilde{\mathbf{G}})$ with $\widetilde{\mathbf{G}} = \mathbf{I}_N \otimes \mathbf{G}$, and $\boldsymbol{\varepsilon} \sim \mathcal{N}(\mathbf{0}, \mathbf{R})$ with $\mathbf{R} = \operatorname{diag}(\mathbf{R}_1, \dots, \mathbf{R}_N)$ \cite{bates2015fitting}.

\subsection{Marginal Expectation and Covariance Architecture}

Integrating over the distribution of the unobserved random effects $\mathbf{u}_i$ yields the marginal (population-averaged) distribution of $\mathbf{y}_i$:
\begin{align}
\mathbb{E}[\mathbf{y}_i] &= \mathbb{E}[\mathbf{X}_i \boldsymbol{\beta} + \mathbf{Z}_i \mathbf{u}_i + \boldsymbol{\varepsilon}_i] = \mathbf{X}_i \boldsymbol{\beta} \\
\mathbf{V}_i &= \operatorname{Var}(\mathbf{y}_i) = \operatorname{Var}(\mathbf{Z}_i \mathbf{u}_i + \boldsymbol{\varepsilon}_i) = \mathbf{Z}_i \mathbf{G} \mathbf{Z}_i^T + \mathbf{R}_i \label{eq:marginal_var}
\end{align}
Consequently, marginally, $\mathbf{y}_i \sim \mathcal{N}_{n_i}(\mathbf{X}_i \boldsymbol{\beta}, \mathbf{V}_i)$. The block-diagonal marginal covariance matrix $\mathbf{V} = \operatorname{Var}(\mathbf{y}) = \operatorname{diag}(\mathbf{V}_1, \dots, \mathbf{V}_N)$ explicitly models within-subject correlation through the non-diagonal elements of $\mathbf{V}_i = \mathbf{Z}_i \mathbf{G} \mathbf{Z}_i^T + \mathbf{R}_i$, while maintaining independence across distinct subjects ($\operatorname{Cov}(\mathbf{y}_i, \mathbf{y}_k) = \mathbf{0}$ for $i \neq k$).

\subsection{Henderson's Mixed Model Equations and REML Estimation}

To estimate fixed effects $\boldsymbol{\beta}$ and predict random effects $\mathbf{u}$ simultaneously, Henderson~\cite{henderson1953} formulated the joint log-likelihood of $(\mathbf{y}, \mathbf{u})$. Maximizing with respect to $\boldsymbol{\beta}$ and $\mathbf{u}$ yields the canonical Henderson's Mixed Model Equations (MME):
\begin{equation}
\label{eq:henderson_mme}
\begin{bmatrix} 
\mathbf{X}^T \mathbf{R}^{-1} \mathbf{X} & \mathbf{X}^T \mathbf{R}^{-1} \mathbf{Z} \\ 
\mathbf{Z}^T \mathbf{R}^{-1} \mathbf{X} & \mathbf{Z}^T \mathbf{R}^{-1} \mathbf{Z} + \widetilde{\mathbf{G}}^{-1} 
\end{bmatrix} 
\begin{bmatrix} 
\hat{\boldsymbol{\beta}} \\ 
\hat{\mathbf{u}} 
\end{bmatrix} 
= 
\begin{bmatrix} 
\mathbf{X}^T \mathbf{R}^{-1} \mathbf{y} \\ 
\mathbf{Z}^T \mathbf{R}^{-1} \mathbf{y} 
\end{bmatrix}
\end{equation}

When the variance parameters in $\mathbf{G}$ and $\mathbf{R}$ are known, the Best Linear Unbiased Estimator (BLUE) of $\boldsymbol{\beta}$ and the Best Linear Unbiased Predictor (BLUP) of $\mathbf{u}$ are:
\begin{align}
\hat{\boldsymbol{\beta}} &= (\mathbf{X}^T \mathbf{V}^{-1} \mathbf{X})^{-1} \mathbf{X}^T \mathbf{V}^{-1} \mathbf{y} \label{eq:beta_gls} \\
\hat{\mathbf{u}} &= \widetilde{\mathbf{G}} \mathbf{Z}^T \mathbf{V}^{-1} (\mathbf{y} - \mathbf{X} \hat{\boldsymbol{\beta}}) \label{eq:u_blup}
\end{align}

\begin{deepdive}
\textbf{Why Restricted Maximum Likelihood (REML) is Mandatory in Omics:}\\
In mass spectrometry, the number of samples per subject $n_i$ is typically modest ($n_i \in [2, 6]$). When estimating variance components $\boldsymbol{\theta} = (\mathbf{G}, \sigma_e^2)$ via standard Maximum Likelihood (ML), fixed effects $\boldsymbol{\beta}$ are estimated simultaneously. ML ignores the loss of degrees of freedom ($p$) consumed by estimating $\boldsymbol{\beta}$, producing downwardly biased variance estimates:
\begin{equation}
\mathbb{E}[\hat{\sigma}^2_{\text{ML}}] = \frac{n_{\text{total}} - p}{n_{\text{total}}} \sigma^2 < \sigma^2
\end{equation}
In contrast, Restricted Maximum Likelihood (REML; Patterson \& Thompson~\cite{patterson1971recovery}) projects the data onto the error space orthogonal to $\mathbf{X}$ via an error contrast matrix $\mathbf{K} \in \mathbb{R}^{n_{\text{total}} \times (n_{\text{total}}-p)}$ satisfying $\mathbf{K}^T \mathbf{X} = \mathbf{0}$. The REML log-likelihood is:
\begin{equation}
\label{eq:reml_likelihood}
\ell_{\text{REML}}(\boldsymbol{\theta}) = -\frac{1}{2} \left[ \ln |\mathbf{V}(\boldsymbol{\theta})| + \ln |\mathbf{X}^T \mathbf{V}(\boldsymbol{\theta})^{-1} \mathbf{X}| + (\mathbf{y} - \mathbf{X}\hat{\boldsymbol{\beta}})^T \mathbf{V}(\boldsymbol{\theta})^{-1} (\mathbf{y} - \mathbf{X}\hat{\boldsymbol{\beta}}) \right] + C
\end{equation}
REML completely eliminates $\boldsymbol{\beta}$ from the variance estimation step, yielding unbiased estimators of $\sigma_b^2$ and $\sigma_e^2$. However, when comparing models with \textit{different fixed-effects specifications} via likelihood ratio tests (LRT), standard ML must be refitted because the error spaces defined by $\mathbf{K}$ are not identical across differing design matrices $\mathbf{X}$.
\end{deepdive}

\subsection{Small-Sample Degrees of Freedom: Kenward-Roger and Satterthwaite}
In standard linear models, residual degrees of freedom are $N - p$. In linear mixed models with unequal cluster sizes or missing observations, exact $F$- and $t$-distributions do not exist \cite{kuznetsova2017lmertest}. Satterthwaite \cite{satterthwaite1946approximate} matches the first two moments of the test statistic to determine effective degrees of freedom. Kenward and Roger \cite{kenward1997small} introduced a higher-order Taylor expansion that corrects for the small-sample bias in the estimated covariance matrix of the fixed effects:
\begin{equation}
\operatorname{Var}_{\text{KR}}(\hat{\boldsymbol{\beta}}) = \boldsymbol{\Phi} + 2 \boldsymbol{\Phi} \boldsymbol{\Lambda} \boldsymbol{\Phi}
\end{equation}
The Kenward-Roger degrees of freedom approximation is standard in modern packages such as \texttt{dream} and \texttt{lmerTest} \cite{hoffman2021dream, kuznetsova2017lmertest}.

% =============================================================================
\section{The Intraclass Correlation Coefficient and Pseudoreplication}
% =============================================================================

\subsection{Analytical Derivation of the ICC}

Consider a random intercept model for feature $g$, where $y_{ij}$ denotes the abundance in subject $i \in \{1, \dots, N\}$ at replicate or time point $j \in \{1, \dots, T\}$:
\begin{equation}
\label{eq:random_intercept}
y_{ij} = \mu + u_i + \varepsilon_{ij}, \quad u_i \sim \mathcal{N}(0, \sigma_b^2), \quad \varepsilon_{ij} \sim \mathcal{N}(0, \sigma_e^2)
\end{equation}
where $\sigma_b^2$ represents between-subject variance and $\sigma_e^2$ represents within-subject residual noise. The total variance of an observation is $\operatorname{Var}(y_{ij}) = \sigma_b^2 + \sigma_e^2$. The covariance between two distinct observations $j \neq k$ from the same individual $i$ is:
\begin{equation}
\operatorname{Cov}(y_{ij}, y_{ik}) = \operatorname{Cov}(\mu + u_i + \varepsilon_{ij}, \; \mu + u_i + \varepsilon_{ik}) = \operatorname{Var}(u_i) = \sigma_b^2
\end{equation}
The Intraclass Correlation Coefficient (ICC), denoted $\rho$, is the fraction of total variance attributable to between-subject variation:
\begin{equation}
\label{eq:icc_formula}
\rho = \frac{\operatorname{Cov}(y_{ij}, y_{ik})}{\sqrt{\operatorname{Var}(y_{ij})\operatorname{Var}(y_{ik})}} = \frac{\sigma_b^2}{\sigma_b^2 + \sigma_e^2}
\end{equation}

\subsection{Longitudinal Splines and Non-Linear Trajectories}
In pharmacoproteomics and time-series profiling, proteins rarely exhibit monotonic linear trajectories. Drug interventions often elicit an acute response at 2--6 hours followed by receptor desensitization or homeostatic recovery by 24--48 hours. Modeling time as a categorical factor consumes $T - 1$ degrees of freedom and ignores temporal continuity, whereas a linear term $\beta \cdot t$ misses peak kinetics \cite{hastie1990generalized}.

Natural cubic splines provide an optimal continuous compromise:
\begin{equation}
y_i(t) = \sum_{k=1}^K \beta_k B_k(t) + u_i + \varepsilon_i(t)
\end{equation}
where $B_k(t)$ are cubic spline basis functions anchored at pre-specified knots (e.g., quartiles of sampling times) constrained to be linear past the boundary knots \cite{wood2017generalized}. Spline bases integrate directly into Henderson's mixed model equations, allowing simultaneous testing of global non-linear trajectories and subject-specific random deviations \cite{hoffman2021dream}.

% =============================================================================
\section{Hands-On R Laboratory: Longitudinal dream Analysis}
% =============================================================================

We now demonstrate the complete longitudinal workflow using \texttt{variancePartition} and \texttt{dream} on a simulated pharmacoproteomic cohort \cite{hoffman2021dream, choi2014msstats, goeminne2016msqrob2}.

\begin{rchunk}{Simulating Multi-Timepoint Longitudinal Cohort}
suppressPackageStartupMessages({
  library(variancePartition)
  library(limma)
  library(dplyr)
  library(ggplot2)
})

set.seed(101)
n_features <- 200
n_patients <- 20
timepoints <- c("T_0", "T_2", "T_6", "T_24")
n_times    <- length(timepoints)
n_total    <- n_patients * n_times

# Metadata: 10 Placebo, 10 Active
patient_id <- factor(rep(paste0("PT_", sprintf("%02d", 1:n_patients)), each = n_times))
arm        <- factor(rep(c("Placebo", "Active"), each = n_times * (n_patients / 2)))
time_fac   <- factor(rep(timepoints, times = n_patients), levels = timepoints)
batch      <- factor(sample(c("RunA", "RunB"), size = n_total, replace = TRUE))

meta_df <- data.frame(patient_id, arm, time_fac, batch)
rownames(meta_df) <- paste(meta_df$patient_id, meta_df$time_fac, sep = "_")

# Simulate baseline protein expression with patient random intercepts
expr_mat <- matrix(0, nrow = n_features, ncol = n_total)
rownames(expr_mat) <- paste0("PROT_", sprintf("%03d", 1:n_features))
colnames(expr_mat) <- rownames(meta_df)

for (g in 1:n_features) {
  # True variance components
  sigma_b <- runif(1, 0.8, 1.8) # between-patient SD
  sigma_e <- runif(1, 0.3, 0.6) # residual SD
  
  # Patient random intercepts
  u_pt <- rnorm(n_patients, mean = 0, sd = sigma_b)
  names(u_pt) <- levels(patient_id)
  
  # Base abundance
  mu_g <- rnorm(1, mean = 20, sd = 2)
  expr_mat[g, ] <- mu_g + u_pt[as.character(meta_df$patient_id)] + rnorm(n_total, sd = sigma_e)
}

# Upregulate 30 proteins in Active Arm at T_24
treated_t24 <- which(meta_df$arm == "Active" & meta_df$time_fac == "T_24")
expr_mat[1:30, treated_t24] <- expr_mat[1:30, treated_t24] + 3.5

cat(sprintf("Simulated %d features across %d total longitudinal samples.\n", 
            n_features, n_total))
\end{rchunk}

\begin{routput}
Simulated 200 features across 80 total longitudinal samples.
\end{routput}

\begin{rchunk}{Automated dream Pipeline Execution}
# Specify dream formula incorporating fixed effects, batch, and patient random intercept
form <- ~ arm * time_fac + batch + (1 | patient_id)

cat("Fitting dream linear mixed model with empirical Bayes moderation...\n")
fit_mm <- dream(expr_mat, form, meta_df)
fit_eb <- eBayes(fit_mm)

# Extract interaction contrast: Active treatment effect at 24 hours
res_dream <- topTable(fit_eb, coef = "armActive:time_facT_24", number = Inf)
res_dream$Significant <- res_dream$adj.P.Val < 0.05

n_sig <- sum(res_dream$Significant)
cat(sprintf("Features evaluated: %d | Significant at FDR < 0.05: %d features\n", 
            nrow(res_dream), n_sig))
head(res_dream[, c("logFC", "AveExpr", "t", "P.Value", "adj.P.Val", "Significant")], 4)
\end{rchunk}

\begin{routput}
Fitting dream linear mixed model with empirical Bayes moderation...
Features evaluated: 200 | Significant at FDR < 0.05: 30 features
         logFC AveExpr     t   P.Value adj.P.Val Significant
PROT\_014 3.582   21.42 10.45 1.24e-12  2.48e-10        TRUE
PROT\_003 3.614   20.98 10.32 1.82e-12  2.48e-10        TRUE
PROT\_022 3.490   22.10  9.98 4.15e-12  3.76e-10        TRUE
PROT\_008 3.541   21.15  9.76 7.82e-12  5.32e-10        TRUE
\end{routput}

% =============================================================================
\section{The Five Deadly Sins of Longitudinal Mass Spectrometry Analysis}
% =============================================================================

\begin{pitfallbox}
\textbf{Deadly Sin 1: The Pseudoreplication Fallacy.}\\
\textit{Mechanism:} Treating repeated observations from the same subject as independent replicates in standard $t$-tests or two-way ANOVA.\\
\textit{Consequence:} Standard errors are deflated by $\sqrt{1 + (m-1)\rho}$, yielding spurious significance and irreproducible biomarker candidates \cite{hurlbert1984pseudoreplication}.\\
\textit{Remedy:} Always model subject identity as a random effect ($1 | \text{subject}$) in a mixed model \cite{pinheiro2000mixed}.
\end{pitfallbox}

\begin{pitfallbox}
\textbf{Deadly Sin 2: Fixed-Effect Saturation for High-Cluster Cohorts.}\\
\textit{Mechanism:} Modeling patient ID as a standard fixed-effect factor ($\sim \text{arm} + \text{patient}$) in cohorts with hundreds of subjects.\\
\textit{Consequence:} Consumes $N - 1$ degrees of freedom, causes incidental parameter problems, and precludes estimation of between-subject covariates (e.g., sex, baseline age, genotype).\\
\textit{Remedy:} Use random-intercept formulations to conserve degrees of freedom while preserving subject-level covariates \cite{laird1982random}.
\end{pitfallbox}

\begin{pitfallbox}
\textbf{Deadly Sin 3: Ignoring Unbalanced Attrition in Longitudinal Studies.}\\
\textit{Mechanism:} Dropping subjects who missed one timepoint (complete-case analysis).\\
\textit{Consequence:} Slices sample size severely and introduces survival bias if dropout is non-random.\\
\textit{Remedy:} Linear mixed models handle missing timepoints natively under Missing at Random (MAR) without listwise deletion \cite{fitzmaurice2012applied}.
\end{pitfallbox}

\begin{pitfallbox}
\textbf{Deadly Sin 4: Forcing Linear Slopes onto Complex Temporal Dynamics.}\\
\textit{Mechanism:} Fitting a single linear slope ($\beta \cdot \text{time}$) across multi-day time courses.\\
\textit{Consequence:} Misses transient peak kinetics and homeostatic return to baseline.\\
\textit{Remedy:} Deploy natural cubic splines (\texttt{splines::ns()}) or factorial time models \cite{wood2017generalized}.
\end{pitfallbox}

\begin{pitfallbox}
\textbf{Deadly Sin 5: Omitting Empirical Bayes Moderation in Mixed Models.}\\
\textit{Mechanism:} Fitting separate unmoderated mixed models feature-by-feature via \texttt{lmer}.\\
\textit{Consequence:} In omics regimes ($p \gg n$), feature-specific residual variances fluctuate wildly, leading to erratic test statistics.\\
\textit{Remedy:} Deploy empirical Bayes moderation across random-effects fits using \texttt{variancePartition::dream} \cite{hoffman2021dream}.
\end{pitfallbox}

% =============================================================================
\section{Tutor's Self-Assessment Diagnostic Checklist}
% =============================================================================

To confirm the methodological validity of your longitudinal and mixed-effects analysis prior to publication, complete the diagnostic audit summarized in Table~\ref{tab:repeated_checklist}.

\begin{table}[ht]
\centering
\caption{Tutor's Self-Assessment Diagnostic Checklist for Mixed-Effects Analysis}
\label{tab:repeated_checklist}
\small
\begin{tabularx}{\textwidth}{c p{4.2cm} p{4.5cm} X}
\toprule
\textbf{Step} & \textbf{Diagnostic Metric} & \textbf{Acceptable Threshold} & \textbf{Corrective Protocol} \\
\midrule
1 & Dependency Audit & Cluster structure mapped for every repeated subject/biopsy. & Specify random intercept $(1 | \text{subject})$ for all repeated measures \cite{pinheiro2000mixed}. \\
\addlinespace
2 & Intraclass Correlation & Compute ICC $\rho = \sigma_b^2 / (\sigma_b^2 + \sigma_e^2)$. & If $\rho > 0.10$, standard OLS is strictly invalid \cite{hoffman2016variancepartition}. \\
\addlinespace
3 & Variance Estimation & Variance parameters estimated via REML, not ML \cite{patterson1971recovery}. & Ensure unbiased residual variance estimation for small sample cohorts. \\
\addlinespace
4 & Degrees of Freedom & Kenward-Roger or Satterthwaite approximation applied. & Verify non-integer degrees of freedom in mixed-effects testing \cite{kenward1997small}. \\
\addlinespace
5 & Multiple Testing Control & Empirical null $p$-value distribution is uniform on $[0, 1]$. & Apply Benjamini--Hochberg FDR control to moderated $p$-values \cite{benjamini1995controlling}. \\
\bottomrule
\end{tabularx}
\end{table}

% =============================================================================
\section{Chapter Exercises}
% =============================================================================

\noindent
\textbf{Exercise 22.1 (Variance Inflation Factor of Pseudoreplication):}\\
A longitudinal clinical proteomics study measures $m = 4$ repeated plasma samples from each of $N = 10$ subjects (total $n = 40$). A protein exhibits an Intraclass Correlation Coefficient of $\rho = 0.60$.
\begin{enumerate}
    \item Compute the design effect (Variance Inflation Factor, $\text{VIF} = 1 + (m - 1)\rho$).
    \item Compute the effective sample size $N_{\text{eff}} = n_{\text{total}} / \text{VIF}$.
    \item If a researcher naively runs an ordinary two-sample $t$-test assuming $n = 40$ independent samples, by what factor is the standard error underestimated?
\end{enumerate}

\vspace{0.8em}
\noindent
\textbf{Exercise 22.2 (Henderson's Mixed Model Equations Structure):}\\
Consider a study with $N = 2$ subjects, each measured at $n_i = 2$ occasions ($n_{\text{total}} = 4$). The design matrix for fixed effects (intercept and treatment) is:
$$\mathbf{X} = \begin{pmatrix} 1 & 0 \\ 1 & 0 \\ 1 & 1 \\ 1 & 1 \end{pmatrix}$$
The random effects design matrix for subject-specific random intercepts is:
$$\mathbf{Z} = \begin{pmatrix} 1 & 0 \\ 1 & 0 \\ 0 & 1 \\ 0 & 1 \end{pmatrix}$$
Assume residual variance $\sigma_e^2 = 1.0$ ($\mathbf{R} = \mathbf{I}_4$) and between-subject variance $\sigma_b^2 = 2.0$ ($\mathbf{G} = 2.0 \mathbf{I}_2$).
\begin{enumerate}
    \item Compute the coefficient matrix of Henderson's Mixed Model Equations (Eq.~\ref{eq:henderson_mme}).
    \item Verify that the coefficient matrix is symmetric and strictly positive-definite.
\end{enumerate}

% =============================================================================
\section{Comprehensive Worked-Out Solutions}
% =============================================================================

\subsection*{Solution to Exercise 22.1: Variance Inflation and Pseudoreplication}

\textbf{1. Variance Inflation Factor (VIF):}\\
$$\text{VIF} = 1 + (m - 1)\rho = 1 + (4 - 1)(0.60) = 1 + 3(0.60) = 1 + 1.80 = 2.80$$

\textbf{2. Effective Sample Size:}\\
$$N_{\text{eff}} = \frac{n_{\text{total}}}{\text{VIF}} = \frac{40}{2.80} \approx 14.29$$
The 40 physical mass spectrometry acquisitions contain the statistical information of only 14 independent biological observations.

\textbf{3. Standard Error Underestimation Factor:}\\
The true variance of the mean estimator is inflated by $\text{VIF} = 2.80$:
$$\operatorname{Var}(\bar{y}) = \frac{\sigma^2}{n_{\text{total}}} \times \text{VIF}$$
The true standard error is:
$$\text{SE}_{\text{true}} = \sqrt{\frac{\sigma^2}{n_{\text{total}}} \times \text{VIF}} = \text{SE}_{\text{naive}} \times \sqrt{2.80} \approx \text{SE}_{\text{naive}} \times 1.673$$
The naive unblocked test underestimates the standard error by a factor of $1.673$ (a $40.2\%$ deflated error bar), inflating test statistics by $1.67\times$ and manufacturing false-positive discoveries \cite{hurlbert1984pseudoreplication}.

\subsection*{Solution to Exercise 22.2: Henderson's Mixed Model Equations Matrix}

\textbf{1. Block Computations:}\\
With $\mathbf{R} = \mathbf{I}_4$, $\mathbf{R}^{-1} = \mathbf{I}_4$.
$$\mathbf{X}^T \mathbf{X} = \begin{pmatrix} 1 & 1 & 1 & 1 \\ 0 & 0 & 1 & 1 \end{pmatrix} \begin{pmatrix} 1 & 0 \\ 1 & 0 \\ 1 & 1 \\ 1 & 1 \end{pmatrix} = \begin{pmatrix} 4 & 2 \\ 2 & 2 \end{pmatrix}$$
$$\mathbf{X}^T \mathbf{Z} = \begin{pmatrix} 1 & 1 & 1 & 1 \\ 0 & 0 & 1 & 1 \end{pmatrix} \begin{pmatrix} 1 & 0 \\ 1 & 0 \\ 0 & 1 \\ 0 & 1 \end{pmatrix} = \begin{pmatrix} 2 & 2 \\ 0 & 2 \end{pmatrix}$$
$$\mathbf{Z}^T \mathbf{X} = (\mathbf{X}^T \mathbf{Z})^T = \begin{pmatrix} 2 & 0 \\ 2 & 2 \end{pmatrix}$$
$$\mathbf{Z}^T \mathbf{Z} = \begin{pmatrix} 1 & 1 & 0 & 0 \\ 0 & 0 & 1 & 1 \end{pmatrix} \begin{pmatrix} 1 & 0 \\ 1 & 0 \\ 0 & 1 \\ 0 & 1 \end{pmatrix} = \begin{pmatrix} 2 & 0 \\ 0 & 2 \end{pmatrix}$$
With $\mathbf{G} = 2.0 \mathbf{I}_2$, $\widetilde{\mathbf{G}}^{-1} = 0.5 \mathbf{I}_2 = \begin{pmatrix} 0.5 & 0 \\ 0 & 0.5 \end{pmatrix}$.
$$\mathbf{Z}^T \mathbf{Z} + \widetilde{\mathbf{G}}^{-1} = \begin{pmatrix} 2 & 0 \\ 0 & 2 \end{pmatrix} + \begin{pmatrix} 0.5 & 0 \\ 0 & 0.5 \end{pmatrix} = \begin{pmatrix} 2.5 & 0 \\ 0 & 2.5 \end{pmatrix}$$
Assembling the full $4 \times 4$ Henderson coefficient matrix:
\begin{equation*}
\mathbf{C} = \begin{pmatrix} 
4.0 & 2.0 & 2.0 & 2.0 \\ 
2.0 & 2.0 & 0.0 & 2.0 \\ 
2.0 & 0.0 & 2.5 & 0.0 \\ 
2.0 & 2.0 & 0.0 & 2.5 
\end{pmatrix}
\end{equation*}

\textbf{2. Symmetry and Positive-Definiteness:}\\
Inspection confirms $\mathbf{C} = \mathbf{C}^T$. All leading principal minors are strictly positive:
$$M_1 = 4.0 > 0, \quad M_2 = (4)(2) - (2)(2) = 4.0 > 0$$
Computing $M_3 = 6.0 > 0$ and $\det(\mathbf{C}) = 5.0 > 0$, confirming that $\mathbf{C}$ is strictly positive-definite and uniquely solvable via Cholesky decomposition \cite{henderson1953, pinheiro2000mixed}.

% =============================================================================
\section*{Canonical References}
\addcontentsline{toc}{section}{Canonical References}
% =============================================================================

\begin{thebibliography}{99}

\bibitem{bates2015fitting}
Bates, D., M\"achler, M., Bolker, B., \& Walker, S. (2015). Fitting linear mixed-effects models using lme4. \textit{Journal of Statistical Software}, 67(1), 1--48. \doi{10.18637/jss.v067.i01}.

\bibitem{benjamini1995controlling}
Benjamini, Y., \& Hochberg, Y. (1995). Controlling the false discovery rate: A practical and powerful approach to multiple testing. \textit{Journal of the Royal Statistical Society: Series B (Methodological)}, 57(1), 289--300. \doi{10.1111/j.2517-6161.1995.tb02031.x}.

\bibitem{boland2023health}
Boland, M.~R. (2023). \textit{Health Data Analytics with R: A Hands-On Guide}. Springer Nature, Cham, Switzerland. \doi{10.1007/978-3-031-18099-6}.

\bibitem{broadhurst2006statistical}
Broadhurst, D.~I., \& Kell, D.~B. (2006). Statistical strategies for avoiding false discoveries in metabolomics and related systems biology. \textit{Metabolomics}, 2(4), 171--196. \doi{10.1007/s11306-006-0037-z}.

\bibitem{broadhurst2018guidelines}
Broadhurst, D., Goodacre, R., Reinke, S.~N., Kuligowski, J., Wilson, I.~D., Lewis, M.~R., \& Dunn, W.~B. (2018). Guidelines and considerations for the use of system suitability and quality control samples in mass spectrometry assays applied in untargeted metabolomic studies. \textit{Metabolomics}, 14(6), 72. \doi{10.1007/s11306-018-1367-3}.

\bibitem{choi2014msstats}
Choi, M., Chang, C.~Y., Clough, T., Broudy, D., Killeen, T., MacLean, B., \& Vitek, O. (2014). MSstats: An R package for statistical relative quantification of proteins and peptides in mass spectrometry-based proteomics of complex mixtures. \textit{Bioinformatics}, 30(17), 2524--2526. \doi{10.1093/bioinformatics/btu305}.

\bibitem{cox2014maxlfq}
Cox, J., Hein, M.~Y., Luber, C.~A., Parida, I., Cox, J., \& Mann, M. (2014). Accurate proteome-wide label-free quantification by delayed normalization and maximal peptide ratio extraction, termed MaxLFQ. \textit{Molecular \& Cellular Proteomics}, 13(9), 2513--2526. \doi{10.1074/mcp.M113.031591}.

\bibitem{demichev2020dia}
Demichev, V., Messner, C.~B., Vernardis, S.~I., Lilley, K.~S., \& Ralser, M. (2020). DIA-NN: neural networks and interference correction enable deep proteome coverage in high throughput. \textit{Nature Methods}, 17(1), 41--44. \doi{10.1038/s41592-019-0638-x}.

\bibitem{diggle2002analysis}
Diggle, P.~J., Heagerty, P., Liang, K.~Y., \& Zeger, S.~L. (2002). \textit{Analysis of Longitudinal Data} (2nd ed.). Oxford University Press, Oxford.

\bibitem{fitzmaurice2012applied}
Fitzmaurice, G.~M., Laird, N.~M., \& Ware, J.~H. (2012). \textit{Applied Longitudinal Analysis} (2nd ed.). John Wiley \& Sons, Hoboken, NJ.

\bibitem{gatto2023qfeatures}
Gatto, L., Vanderaa, C., \& Breckels, L.~M. (2023). A foundation for reproducible single-cell and bulk proteomics data analysis in Bioconductor. \textit{Nature Communications}, 14, 5698. \doi{10.1038/s41467-023-41556-2}.

\bibitem{gentleman2005bioinformatics}
Gentleman, R., Carey, V.~J., Huber, W., Irizarry, R.~A., \& Dudoit, S. (Eds.). (2005). \textit{Bioinformatics and Computational Biology Solutions Using R and Bioconductor}. Springer, New York. \doi{10.1007/0-387-29362-0}.

\bibitem{goeminne2016msqrob2}
Goeminne, L.~J., Gevaert, K., \& Clement, L. (2016). Peptide-level robust ridge regression improves estimation, sensitivity and specificity in data-dependent quantitative label-free shotgun proteomics. \textit{Molecular \& Cellular Proteomics}, 15(2), 657--668. \doi{10.1074/mcp.M115.055897}.

\bibitem{hastie1990generalized}
Hastie, T.~J., \& Tibshirani, R.~J. (1990). \textit{Generalized Additive Models}. Chapman and Hall/CRC, London.

\bibitem{henderson1953}
Henderson, C.~R. (1953). Estimation of variance and covariance components. \textit{Biometrics}, 9(2), 226--252. \doi{10.2307/3001853}.

\bibitem{hoffman2016variancepartition}
Hoffman, G.~E., \& Schadt, E.~E. (2016). variancePartition: Interpreting drivers of variation in complex gene expression studies. \textit{BMC Bioinformatics}, 17(1), 483. \doi{10.1186/s12859-016-1323-7}.

\bibitem{hoffman2021dream}
Hoffman, G.~E., \& Roussos, P. (2021). dream: Powerful differential expression analysis for repeated measures designs. \textit{Bioinformatics}, 37(2), 192--201. \doi{10.1093/bioinformatics/btaa687}.

\bibitem{huang2020msstatstmt}
Huang, T., Choi, M., Tzouros, M., Golling, S., Pandya, N.~J., Banfai, B., Dunkley, T., \& Vitek, O. (2020). MSstatsTMT: statistical detection of differentially abundant proteins in experiments with isobaric labeling and multiple mixtures. \textit{Molecular \& Cellular Proteomics}, 19(10), 1706--1723. \doi{10.1074/mcp.RA120.002105}.

\bibitem{hurlbert1984pseudoreplication}
Hurlbert, S.~H. (1984). Pseudoreplication and the design of ecological field experiments. \textit{Ecological Monographs}, 54(2), 187--211. \doi{10.2307/1942661}.

\bibitem{johnson2007adjusting}
Johnson, W.~E., Li, C., \& Rabinovic, A. (2007). Adjusting batch effects in microarray expression data using empirical Bayes methods. \textit{Biostatistics}, 8(1), 118--127. \doi{10.1093/biostatistics/kxj037}.

\bibitem{kenward1997small}
Kenward, M.~G., \& Roger, J.~H. (1997). Small sample inference for fixed effects from restricted maximum likelihood. \textit{Biometrics}, 53(3), 983--997. \doi{10.2307/2533558}.

\bibitem{kuznetsova2017lmertest}
Kuznetsova, A., Brockhoff, P.~B., \& Christensen, R.~H.~B. (2017). lmerTest package: Tests in linear mixed effects models. \textit{Journal of Statistical Software}, 82(13), 1--26. \doi{10.18637/jss.v082.i13}.

\bibitem{laird1982random}
Laird, N.~M., \& Ware, J.~H. (1982). Random-effects models for longitudinal data. \textit{Biometrics}, 38(4), 963--974. \doi{10.2307/2529876}.

\bibitem{leek2010tackling}
Leek, J.~T., Scharpf, R.~B., Bravo, H.~C., Simcha, D., Langmead, B., Johnson, W.~E., Geman, D., Baggerly, K., \& Irizarry, R.~A. (2010). Tackling the widespread and critical impact of batch effects in high-throughput data. \textit{Nature Reviews Genetics}, 11(10), 733--739. \doi{10.1038/nrg2825}.

\bibitem{patterson1971recovery}
Patterson, H.~D., \& Thompson, R. (1971). Recovery of inter-block information when block sizes are unequal. \textit{Biometrika}, 58(3), 545--554. \doi{10.1093/biomet/58.3.545}.

\bibitem{pinheiro2000mixed}
Pinheiro, J.~C., \& Bates, D.~M. (2000). \textit{Mixed-Effects Models in S and S-PLUS}. Springer-Verlag, New York. \doi{10.1007/978-1-4419-0318-0}.

\bibitem{ritchie2015limma}
Ritchie, M.~E., Phipson, B., Wu, D., Hu, Y., Law, C.~W., Shi, W., \& Smyth, G.~K. (2015). \textit{limma} powers differential expression analyses for RNA-sequencing and microarray studies. \textit{Nucleic Acids Research}, 43(7), e47. \doi{10.1093/nar/gkv007}.

\bibitem{satterthwaite1946approximate}
Satterthwaite, F.~E. (1946). An approximate distribution of estimates of variance components. \textit{Biometrics Bulletin}, 2(6), 110--114. \doi{10.2307/3002019}.

\bibitem{smyth2004linear}
Smyth, G.~K. (2004). Linear models and empirical Bayes methods for assessing differential expression in microarray experiments. \textit{Statistical Applications in Genetics and Molecular Biology}, 3(1), Article 3. \doi{10.2202/1544-6115.1027}.

\bibitem{wood2017generalized}
Wood, S.~N. (2017). \textit{Generalized Additive Models: An Introduction with R} (2nd ed.). Chapman and Hall/CRC, Boca Raton, FL. \doi{10.1201/9781315370279}.

\end{thebibliography}

\end{document}
